This supplement accompanies the Experiments section. It preserves the
complete supporting tables and separates reported measurements from
protocol details and unresolved run associations. Results from the main
comparisons, component study, alternative updates, and parameter sweeps
retain their own reference runs; their different CAST values are not
pooled.

\section{Simulation Protocols}
\label{supp:protocols}
Matched DICE-RL and CAST comparisons start from the same pretrained
checkpoint and retain the residual architecture, critic-training procedure,
replay recipe, and stated interaction budget. DICE-RL's actor maximizes
critic value with selectively applied BC regularization. CAST replaces
that actor objective with action-target regression. The comparison is
between complete actor updates, including their weighting of replay
states. In online refinement, resulting trajectories and learned critic
weights need not remain identical.

\textbf{Replication and evaluation.}
Robomimic evaluations use 300 episodes. Drifting can and square results
average the last three checkpoints separately for training seeds 42--44;
the main table reports the range of these seed-level averages. Two
DICE-RL can runs end before the common 20K-update budget. LIBERO main
results use 100 rollouts for each of evaluation seeds 10000--10002.
These quantify evaluation variation and are not three independently
trained policies. Both the drifting and flow-matching multitask bases
are first trained on all 90 LIBERO tasks. Refinement then targets the same
eight challenging tasks for both bases. Each selected task has its own
residual and critic, deployed through task-ID routing over the shared
frozen multitask base. Reported means cover this selected subset.
The matched LIBERO recipe uses value-guided selection among ten candidates.
DMC returns average ten episodes of 1000 steps each. Swing-up averages
two training seeds; the remaining reported DMC rows each use one.

\textbf{Actor settings.}
The shared LIBERO recipe uses $K=16$ candidates, weak restoration
$\eta_0=0.05$, radial restoration $\eta_{\mathrm{anc}}=1$, an RMS radius
$\rho=0.05$, normalized critic-field cap $\delta_Q=1$, and combined
step cap $\delta_{\mathrm{tot}}=0.15$. The two main-table settings use
$\eta_Q=1$ and $2$ across all eight tasks, respectively. This parameter
scales actor guidance; it is not the critic optimizer learning rate.
Robomimic and default DMC gradient recipes use $\eta_Q=0.5$ and
$\rho=0$. The DMC quadruped result is an explicitly tuned exception.

\section{Physical Refinement Protocol}
\label{supp:physical}
\textbf{Platform and policy.}
The OpenArm setup operates the seven-joint left follower arm and its
parallel gripper. Chest RealSense D435 and wrist D405 cameras provide RGB
observations, resized to $128\times128$ and encoded by a shared
ResNet-18. Follower joint/gripper state and a task-language embedding
condition a flow-matching policy with a diffusion-transformer backbone.
Ten Euler steps generate 30-step chunks of absolute joint/gripper targets.
Commands are issued at 15 Hz; a new chunk is requested every eight control
steps. The common execution interface limits arm-command changes to
0.6 rad/s and ends attempts after 20 s.

\textbf{Task-specific data.}
Policy pretraining uses 40 placement demonstrations and approximately
60 stacking demonstrations. BC refinement uses 30 additional expert
demonstrations for each task, separate from the pretraining data. For RL
refinement, we collect 40 physical robot rollouts per task and retain all
outcomes, including both successes and failures. DICE-RL and CAST share
the same task-specific rollout pool and outcome labels. Refinement is
performed offline after data collection. These counts are training
episodes; the 20 autonomous evaluation attempts per method and task are
separate.

\textbf{Training.}
BC fine-tunes the base and encoder for 40 additional epochs; the documented
peak learning rates are $3\times10^{-5}$ for placement and $10^{-5}$ for
stacking. Both RL arms freeze the base and encoder and train a residual
and five-critic ensemble for 6000 critic updates, with 500 critic-only
warm-up steps and an actor update every two steps. They share the data,
initialization seed, replay-index schedule, and optimizer budget. The
batch size is 32, the actor uses 16 noise samples, and critic hidden
layers have width 256. Actor and
critic learning rates are $10^{-4}$, the target-critic Polyak coefficient
is 0.01, the discount is 0.99, and the CQL-style penalty weight is 0.1.
DICE-RL uses the residual/Q-filtered BC objective with BC weight 100.
CAST uses gradient guidance with $\eta_Q=1$, $\eta_0=0.05$,
$\eta_{\mathrm{anc}}=1$, $\rho=0.05$ RMS, $\delta_Q=1$, and
$\delta_{\mathrm{tot}}=0.15$. These are offline physical-refinement
settings, not a reproduction of every published online DICE-RL setting.
BC versus RL changes data, supervision, and trainable parameters; it is
not a matched collection-effort or actor-only comparison.

\textbf{Replay and evaluation.}
Saved follower state and issued commands provide the state/action pairs
for autonomous attempts. The critic and residual correction operate on
the nominal eight-action executed prefix of each 30-action chunk.
Behavior latents were not recorded, so the critic is noise-independent;
new actor candidates still use same-noise base/residual pairing. Only
successful terminal transitions receive reward 1, all other rewards are
zero, and nonterminal critic targets bootstrap future value estimates.
The autonomous success counts are separate from training losses and
held-out imitation errors. There are 20 attempts per task and method;
placement poses are randomized on the stairs. Paired resets and a common
post-release observation interval are not documented. Historical camera
recordings lack per-frame timestamps, so offline image/state alignment
is approximate. The reported evaluation counts do not yet identify the
exact evaluated checkpoint paths or hashes. Matching the training protocol
therefore does not by itself verify the evaluated checkpoints.

\section{Complete Robomimic Results}
\label{supp:robomimic}
% Sources: tab:matched and the paired/external blocks of tab:baselines,
% iclr2026/sections/05_experiments.tex at 800cfe6. See RESULTS_PROVENANCE.md.
\begin{table}[H]
\centering
\caption{Robomimic success (\%). Top: DICE-RL and CAST refine the same
pretrained base within each row. Drifting entries span three training
seeds, each averaged over its last three 300-episode evaluations;
$\dagger$ marks two DICE-RL can runs ending before 20K updates.
Diffusion results use 496K environment steps; $z$ denotes value-only
CAST with an unspecified interaction budget. Bottom: external square
results use separate bases, budgets, and evaluation protocols.
$\Delta$ is the gain over each method's own base in percentage points
(pp); $*$ marks a stricter success rule and $\circ$ an ongoing run.
Bold marks the highest refinement result per row in the top block,
and the highest fine-tuned result and gain within each external group.}
\label{tab:robomimic}
\small
\setlength{\tabcolsep}{3pt}
\begin{tabular}{@{}lccc@{}}
\toprule
Task / base class & Base & DICE-RL & CAST \\
\midrule
Can / drifting & 87.7 & 94.0--97.0$^{\dagger}$ & \textbf{97.8--99.3} \\
Square / drifting & 38.2 & 86.1--90.9 & \textbf{91.2--93.0} \\
Square / diffusion & 60.0 & \textbf{96.5} & 95.9 \\
Square / flow matching & 57.0 & 83.1 & \textbf{93.4}$^{z}$ \\
\midrule
\multicolumn{4}{@{}l}{\emph{External square references (separate protocols)}} \\
Method & Own base & Fine-tuned & $\Delta$ (pp) \\
\midrule
DIPO~\cite{dipo2024}$^{*}$ & 34.8 & 15.6 & $-19.2$ \\
QSM~\cite{psenka2024qsm}$^{*}$ & 34.8 & 0.0 & $-34.8$ \\
DQL$^{*}$ & 34.8 & 0.0 & $-34.8$ \\
IDQL~\cite{hansenestruch2023idql}$^{*}$ & 34.8 & 32.7 & $-2.1$ \\
AWR~\cite{peng2019awr}$^{*}$ & 34.8 & \textbf{39.3} & $\mathbf{+4.5}$ \\
DPPO~\cite{ren2024dppo} & 60.0 & \textbf{67.8}$^{\circ}$ & $\mathbf{+7.8}$ \\
\bottomrule
\end{tabular}
\end{table}

Table~\ref{tab:robomimic} retains all reported generator-class comparisons.
On drifting square, CAST is ahead at the 20K-update comparison, while
both actors approach 97--99\% with longer training near 40K updates.
The flow-matching square entry uses value-only guidance; its exact
interaction budget is unspecified and no speed claim is drawn from it.

The diffusion square comparison uses a shared base with 60\% success.
The three first evaluated checkpoints at or above 90\% success are
128K, 160K, and 160K steps for CAST, and 224K, 240K, and 256K for
DICE-RL. Their means are $448/3\approx149.3$K and 240K steps, a
$1-(448/3)/240\approx37.8\%$ reduction. These are discrete evaluation
checkpoints, not exact crossing times interpolated from a learning curve.
At 496K steps, CAST reaches 95.9\% and DICE-RL reaches 96.5\%.
No confidence interval or continuous trajectory is inferred from these
summary measurements.

\textbf{Separate-protocol external references.}
DIPO, QSM, DQL, IDQL, and AWR use implementations adapted by the DPPO
authors and approximately 24M environment steps. Their stricter rule
requires success across an action chunk and scores the base at 34.8\%.
All gains in the external block use that same-harness base: for example,
AWR's 39.3\% is a 4.5 pp improvement. DPPO's 67.8\% comes from an
ongoing run on a 60.0\% base; a separate original-architecture run
reaches 88.7\% at 11.7M steps. Differences in bases, systems, budgets,
and evaluation prevent interpreting these rows as an isolated test of
CAST's components.

\section{Component and Design Studies}
\label{supp:components}
% Source: ICLR tab:constraint; retain this campaign's CAST reference.
% Launcher B_ANC omits the weak pull; this is not an exact 2x2 factorial.
\begin{table}[H]
\centering
\caption{LIBERO success (\%) with different restoration and clipping
controls. Neither disables all controls; Clip enables only the two caps;
Restore enables only radial restoration ($\eta_0=0$); CAST enables both
restoring terms and both caps. Restore--CAST therefore changes weak
restoration as well as clipping. Entries are task aggregates without
training-variance estimates; Mean averages the five tasks. Bold marks
the highest success per row, including the base.}
\label{tab:ablations}
\small
\setlength{\tabcolsep}{3pt}
\begin{tabular}{@{}lccccc@{}}
\toprule
Task & Base & Neither & Clip & Restore & CAST \\
\midrule
Red mug & 57.3 & 0.0 & 0.0 & \textbf{81.0} & 60.0 \\
Ketchup & 61.0 & 0.0 & 0.0 & 38.7 & \textbf{62.7} \\
Bowl & 81.8 & 0.0 & 0.0 & 80.7 & \textbf{84.0} \\
Juice & 78.5 & 0.0 & 0.0 & \textbf{86.0} & 82.0 \\
Right caddy & 57.8 & 0.0 & 0.0 & 42.3 & \textbf{63.0} \\
\midrule
Mean & 67.3 & 0.0 & 0.0 & 65.7 & \textbf{70.3} \\
\bottomrule
\end{tabular}
\end{table}

\textbf{Component settings and replication.}
The four configurations are specified in Table~\ref{supp:control-settings};
Table~\ref{tab:ablations} preserves their exact reported success rates.
The source manuscript describes three seeds per cell, but the available
aggregate record does not contain per-seed measurements or run manifests
that verify training/evaluation seed linkage for every cell. We therefore
show task aggregates without training-variance estimates. The main chart's
points are tasks, not seeds. Its means are computed from the five task
values before rounding: base 67.28\%, radial restoration 65.74\%, and
full CAST 70.34\%. Their mean changes from the base are $-1.54$ and
$+3.06$ pp, respectively.

\begin{table}[H]
\centering
\caption{Controls enabled in the LIBERO component comparison. Caps
limit the critic field and the combined displacement. Radial restoration
only also disables weak restoration, so the study is not an exact
two-factor ablation.}
\label{supp:control-settings}
\small
\begin{tabular}{@{}lccc@{}}
\toprule
Configuration & Weak restoration & Radial restoration & Caps \\
\midrule
No controls & Off & Off & Off \\
Clipping only & Off & Off & On \\
Radial restoration only & Off & On & Off \\
Full CAST & On & On & On \\
\bottomrule
\end{tabular}
\end{table}

\textbf{Alternative constraints.}
The loss-space hinge replaces the residual penalty's shape with a penalty
on excess RMS departure inside DICE-RL's actor objective. It retains
DICE-RL's filtering and weighting. The LIBERO hard-projection recipe
removes both restoring terms and both displacement caps, then projects
the action target into a base-centered RMS ball. Thus, the alternatives
change complete update rules rather than constraint location alone.
The source reports three training seeds, with 300-episode square
evaluations and 100 rollouts for each of three LIBERO evaluation seeds.
The constraint campaign uses CAST references of 97/78/68\% for
square/mug/ketchup, distinct from the guidance-study references.
The square value is the experimenter's corrected measurement for this
comparison; the other campaigns retain their reported values.

The LIBERO launchers identify the shared drifting multitask base, for
which the source reports 57.3\% on mug and 61.0\% on ketchup. Both square
design studies use the drifting base with 38.2\% success. The constraint
comparison reports 48.0\% for the loss-space hinge and 97.0\% for CAST.
The guidance study retains its own CAST reference and auxiliary settings.
The separate
diffusion-base efficiency study still starts from 60.0\% success.

\textbf{Alternative guidance.}
The gradient recipe uses normalized critic action derivatives. Top-$k$
transport uses the highest-valued current candidates as attractors;
value-tilted transport assigns weights proportional to
$\exp(\hat q_i/\tau)$, where $\hat q_i$ is standardized within the
state's candidate set. Both value-only recipes use kernel-weighted
attraction and repulsion over sampled actions. In the available launchers,
they also use a weak distributional pull toward independent base samples,
whereas the gradient recipe uses the paired pointwise weak pull. Historical
square recipes additionally differ in radial-restoration settings. The
reported square ranges span three training seeds. These comparisons do
not isolate the source of critic guidance or establish that gradients
always outperform rankings.

\textbf{Restoration radius.}
Table~\ref{supp:radius} gives the exploratory single-training-seed sweep
cited in the main text. The radius is measured in normalized RMS action
units and controls when stronger restoration starts; it is not a hard
bound on fitted policy outputs. The weak pull remains at infinite radius.
These sweep references are retained separately from the main and
component-study CAST results.
\begin{table}[H]
\centering
\caption{LIBERO success (\%) across restoration radii, with one training
seed per setting. Radius $\rho$ is measured in normalized RMS action
units; $\rho=0.05$ is this campaign's reference. Infinite radius disables
radial restoration while retaining the weak pull. These results do not
establish an optimal radius.}
\label{supp:radius}
\small
\begin{tabular}{@{}lrr@{}}
\toprule
Radius $\rho$ (RMS) & Mug & Ketchup \\
\midrule
0 & 54.9 & 60.1 \\
0.02 & 86.3 & 73.4 \\
0.05 & 78.1 & 68.4 \\
0.10 & 81.2 & 69.4 \\
0.20 & 77.1 & 67.1 \\
$\infty$ & 80.3 & 67.1 \\
\bottomrule
\end{tabular}
\end{table}

\section{Broader Evaluations}
\label{supp:scope}
% Source: tab:dmc-scope, including every incomplete-run marker.
\begin{table}[H]
\centering
\caption{DMC returns for refinements of a shared frozen base per task.
Each evaluation averages ten 1000-step episodes (maximum return: 1000).
The target budget is 400K environment steps; $\circ$ marks a run below
budget, $v$ high evaluation variance, and $t$ a tuned critic-guidance
scale for the actor. Swing-up averages two training seeds; other tasks
use one. Bold marks the highest reported return per row, including ties
and runs below budget.}
\label{tab:dmc}
\small
\setlength{\tabcolsep}{3.5pt}
\begin{tabular}{@{}lrrrr@{}}
\toprule
Task & Base & Unconstrained & DICE-RL & CAST \\
\midrule
\multicolumn{5}{@{}l}{\emph{Dense reward}} \\
Cartpole swing-up & 429 & 575 & 318 & \textbf{715} \\
Cartpole balance & 378 & 815$^{\circ}$ & 844 & \textbf{911} \\
Quadruped run & 238 & 97$^{\circ}$ & 283 & \textbf{312}$^{\circ t}$ \\
Finger spin & 3 & 0$^{\circ}$ & \textbf{187}$^{\circ}$ & \textbf{187}$^{\circ}$ \\
\midrule
\multicolumn{5}{@{}l}{\emph{Sparse reward}} \\
Cartpole balance & 90 & 880$^{\circ}$ & 843$^{v}$ & \textbf{1000} \\
Ball-in-cup catch & 86 & \textbf{940}$^{\circ}$ & 66$^{\circ}$ & 925 \\
\bottomrule
\end{tabular}
\end{table}

DMC evaluates one-step generative bases distilled from SAC experience.
Table~\ref{tab:dmc} retains dense/sparse groupings and all incomplete-run,
tuning, and high-variance markers. Returns use the original 0--1000 scale.
The target budget is 400K environment steps; exact budgets for the marked
interim checkpoints are not available in the local records. Quadruped's
CAST return of 312 uses $\eta_Q=1$ and is below budget; its default
$\eta_Q=0.5$ arm reports 267. The table is not a ranking of all methods
at completed, matched budgets. In particular, the unconstrained sparse
ball-in-cup interim return of 940 exceeds CAST's 925. Sparse reward alone
therefore does not imply failure without restoration.

\begin{table}[H]
\centering
\caption{Success (\%) in two comparisons where DICE-RL exceeds CAST.
Robomimic transport uses matched budgets of 2.24M environment steps.
Shared-critic LIBERO jointly refines over eight tasks and evaluates
best-of-16 value-guided selection; this compares the complete systems
under their respective refinement recipes.}
\label{supp:scope-reversals}
\small
\begin{tabular}{@{}lrrl@{}}
\toprule
Study & DICE-RL & CAST & Protocol \\
\midrule
Robomimic transport & 91.2 & 59.4 & 2.24M environment steps \\
Shared-critic LIBERO & 72.2 & 69.2 & Best-of-16; base 66.9\% \\
\bottomrule
\end{tabular}
\end{table}

Transport shows that restoration can accompany worse refinement even at
a matched budget; it does not measure how far the optimal behavior lies
from the base. Shared-critic LIBERO jointly trains over the eight tasks,
unlike the main per-task study. The DICE-RL row uses its published recipe,
so this is a system comparison. Single-candidate evaluation and
best-of-16 selection give different conclusions: useful candidate ranking
can remain even when residual learning changes little. Neither comparison
establishes critic accuracy or a universal preference for stronger
restoration.
